\documentclass[11pt,a4paper]{article}

\usepackage[utf8]{inputenc}
\usepackage[T1]{fontenc}
\usepackage[brazilian]{babel}
\usepackage{graphicx}
\usepackage{geometry}
\usepackage{array}
\usepackage{enumitem}
\usepackage{fancyhdr}
\usepackage{lastpage}
\usepackage{multirow}
\usepackage{listingsutf8}
\usepackage{xcolor}
\usepackage{titlesec}
\usepackage{booktabs}
\usepackage{tcolorbox}
\usepackage{hyperref}

\tcbuselibrary{skins,breakable}

\geometry{
  top=2.0cm,
  bottom=2.0cm,
  left=2.5cm,
  right=2.0cm
}

\setlength{\parindent}{0pt}
\setlength{\parskip}{0.4em}
\setlength{\headheight}{14.5pt}
\renewcommand{\arraystretch}{1.35}
\setlist[itemize]{itemsep=0.1em, topsep=0.2em, parsep=0pt, partopsep=0pt}
\setlist[enumerate]{itemsep=0.15em, topsep=0.2em, parsep=0pt, partopsep=0pt}

\pagestyle{fancy}
\fancyhf{}
\lhead{Administração de Redes -- 2026}
\rhead{Domínios Virtuais com NGINX e DNS}
\cfoot{Página \thepage\ de \pageref{LastPage}}

% ── Cores ─────────────────────────────────────────────────────────────────────
\definecolor{codebg}{RGB}{248,248,248}
\definecolor{codeframe}{RGB}{210,210,210}
\definecolor{myblue}{RGB}{0,60,130}
\definecolor{mygreen}{RGB}{30,130,30}
\definecolor{mygray}{RGB}{120,120,120}
\definecolor{mymauve}{RGB}{140,0,140}
\definecolor{infobg}{RGB}{230,242,255}
\definecolor{alertbg}{RGB}{255,248,220}
\definecolor{successbg}{RGB}{230,255,230}

% ── Estilos de código ──────────────────────────────────────────────────────────
\lstdefinestyle{bash}{
  backgroundcolor=\color{codebg},
  frame=single,
  rulecolor=\color{codeframe},
  basicstyle=\ttfamily\footnotesize,
  breaklines=true,
  columns=flexible,
  keepspaces=true,
  showstringspaces=false,
  morecomment=[l]{\#},
  commentstyle=\color{mygray}\itshape,
  stringstyle=\color{mymauve},
  xleftmargin=4pt,
  xrightmargin=4pt,
  aboveskip=4pt,
  belowskip=4pt,
}

\lstdefinestyle{conf}{
  backgroundcolor=\color{codebg},
  frame=single,
  rulecolor=\color{codeframe},
  basicstyle=\ttfamily\footnotesize,
  breaklines=true,
  columns=flexible,
  keepspaces=true,
  showstringspaces=false,
  morecomment=[l]{\#},
  morecomment=[l]{;},
  commentstyle=\color{mygray}\itshape,
  xleftmargin=4pt,
  xrightmargin=4pt,
  aboveskip=4pt,
  belowskip=4pt,
}

\lstdefinestyle{html}{
  backgroundcolor=\color{codebg},
  frame=single,
  rulecolor=\color{codeframe},
  basicstyle=\ttfamily\footnotesize,
  breaklines=true,
  columns=flexible,
  keepspaces=true,
  showstringspaces=false,
  morestring=[b]{"},
  stringstyle=\color{mymauve},
  morecomment=[s]{<!--}{-->},
  commentstyle=\color{mygray}\itshape,
  morekeywords={DOCTYPE,html,head,body,meta,title,h1,h2,p,strong,em,lang,charset,style},
  keywordstyle=\color{myblue},
  xleftmargin=4pt,
  xrightmargin=4pt,
  aboveskip=4pt,
  belowskip=4pt,
}

% ── Caixas coloridas ───────────────────────────────────────────────────────────
\newtcolorbox{infobox}{
  colback=infobg, colframe=myblue!50, boxrule=0.5pt,
  arc=3pt, left=6pt, right=6pt, top=4pt, bottom=4pt,
  breakable
}
\newtcolorbox{alertbox}{
  colback=alertbg, colframe=orange!70, boxrule=0.5pt,
  arc=3pt, left=6pt, right=6pt, top=4pt, bottom=4pt,
  breakable
}
\newtcolorbox{successbox}{
  colback=successbg, colframe=mygreen!60, boxrule=0.5pt,
  arc=3pt, left=6pt, right=6pt, top=4pt, bottom=4pt,
  breakable
}

% ── Seções ─────────────────────────────────────────────────────────────────────
\titleformat{\section}
  {\normalfont\large\bfseries\color{myblue}}
  {\thesection.}{0.5em}{}[\vspace{-0.3em}\color{myblue!40}\rule{\linewidth}{0.5pt}]

\titleformat{\subsection}
  {\normalfont\normalsize\bfseries}
  {\thesubsection}{0.5em}{}

\titlespacing*{\section}{0pt}{1.0em}{0.5em}
\titlespacing*{\subsection}{0pt}{0.7em}{0.3em}

% ─────────────────────────────────────────────────────────────────────────────
\begin{document}

% ── Cabeçalho institucional ────────────────────────────────────────────────────
\begin{tabular}{|>{\centering\arraybackslash}m{3.2cm}|>{\centering\arraybackslash}m{12.0cm}|}
  \hline
  \includegraphics[height=2.5cm]{logo.png}
  &
  \begin{minipage}[c]{12.0cm}
    \centering
    {\large \textbf{Curso de Sistemas de Informação}}\\[0.15em]
    {\large \textbf{Disciplina: Administração de Redes}}\\[0.15em]
    {\large \textbf{Atividade Prática -- Domínios Virtuais com NGINX e DNS}}
  \end{minipage}
  \\
  \hline
\end{tabular}

\vspace{0.6em}

\begin{center}
  \begin{tabular}{|>{\bfseries\arraybackslash}m{3.5cm}|>{\arraybackslash}m{10.2cm}|}
    \hline
    Professor      & Rafael Silva Guimarães \\
    \hline
    Prazo de entrega & 09 de junho de 2026 \\
    \hline
    Valor          & 10,0 pontos \\
    \hline
    Modalidade     & Trabalho em dupla \\
    \hline
  \end{tabular}
\end{center}

\vspace{0.3em}

% ── Descrição ──────────────────────────────────────────────────────────────────
\begin{tcolorbox}[colback=myblue!5, colframe=myblue, boxrule=0.7pt, arc=3pt,
  title={\textbf{Descrição da Atividade}}, fonttitle=\bfseries\color{white},
  colbacktitle=myblue]

O trabalho consiste em publicar \textbf{seis sites estáticos} no NGINX com \textbf{domínios virtuais}, onde a resolução de nomes é feita pelo servidor DNS \textbf{BIND (named)}, configurado via \textbf{Webmin}.

Cada dupla usa os primeiros nomes dos integrantes para compor os domínios:

\vspace{0.3em}
\begin{itemize}
  \item \textbf{Domínios:} \texttt{nome1oaluno.info} \ e \ \texttt{nome2oaluno.info}
  \item \textbf{Hosts} (repetidos em cada domínio): \texttt{marataizes}, \texttt{castelo}, \texttt{cachoeiro}
\end{itemize}

\vspace{0.3em}
\begin{center}
\begin{tabular}{|l|l|}
  \hline
  \textbf{Domínio 1} & \textbf{Domínio 2} \\
  \hline
  \texttt{marataizes.nome1oaluno.info} & \texttt{marataizes.nome2oaluno.info} \\
  \texttt{castelo.nome1oaluno.info}    & \texttt{castelo.nome2oaluno.info} \\
  \texttt{cachoeiro.nome1oaluno.info}  & \texttt{cachoeiro.nome2oaluno.info} \\
  \hline
\end{tabular}
\end{center}

\vspace{0.3em}
Cada um dos seis endereços deve exibir uma \textbf{página HTML própria e distinta}. A entrega é um \textbf{tutorial em PDF} demonstrando as configurações realizadas e os testes de acesso.
\end{tcolorbox}

\vspace{0.3em}

\begin{infobox}
  \textbf{Exemplo utilizado neste tutorial:} \ \texttt{joao.info} (aluno 1) e \texttt{maria.info} (aluno 2), com servidor DNS/Web no IP \texttt{<IP\_VM>}. Substitua pelos nomes reais e pelo IP do ambiente de laboratório.
\end{infobox}

% ── Seção 1 ───────────────────────────────────────────────────────────────────
\section{Ambiente de Configuração}

O ambiente de laboratório utiliza dois equipamentos na mesma rede:

\begin{center}
  \begin{tabular}{|l|l|c|}
    \hline
    \textbf{Máquina} & \textbf{Função} & \textbf{IP (exemplo)} \\
    \hline
    \texttt{servidor} & DNS (BIND9) + Web (NGINX) + Webmin & \texttt{<IP\_HTTP>} \\
    \texttt{cliente}  & Testes de resolução e acesso HTTP   & \texttt{IP\_CLIENTE} \\
    \hline
  \end{tabular}
\end{center}

O Webmin deve estar previamente instalado e acessível em \texttt{https://<IP\_VM>:10000}.

% ── Seção 2 ───────────────────────────────────────────────────────────────────
\section{Configuração do DNS (BIND9) via Webmin}

\subsection{Criando a zona do primeiro domínio}

\begin{enumerate}
  \item Acesse o Webmin: \texttt{https://<IP\_VM>:10000}
  \item No menu lateral: \textbf{Servers} $\rightarrow$ \textbf{BIND DNS Server}.
  \item Clique em \textbf{Create Master Zone} e preencha:
  \begin{itemize}
    \item \textbf{Domain name/Network:} \texttt{joao.info}
    \item \textbf{Email address:} \texttt{admin@joao.info}
    \item \textbf{IP address for zone:} \texttt{ENDERECO\_DNS}
  \end{itemize}
  \item Clique em \textbf{Create}.
\end{enumerate}

\subsection{Adicionando os registros A}

Dentro da zona \texttt{joao.info}, clique em \textbf{Address} e adicione três registros:

\begin{center}
  \begin{tabular}{|l|l|}
    \hline
    \textbf{Name}      & \textbf{Address} \\
    \hline
    \texttt{marataizes} & \texttt{<IP\_VM>} \\
    \texttt{castelo}    & \texttt{<IP\_VM>} \\
    \texttt{cachoeiro}  & \texttt{<IP\_VM>} \\
    \hline
  \end{tabular}
\end{center}

Para cada registro, clique em \textbf{Create}. Ao terminar, clique em \textbf{Apply Configuration}.

\subsection{Criando a zona do segundo domínio}

Repita os passos das seções 2.1 e 2.2 para o segundo domínio (\texttt{maria.info}),
adicionando os mesmos três registros A apontando para \texttt{<IP\_VM>}.

\subsection{Verificando os arquivos de zona gerados}

% ── Seção 3 ───────────────────────────────────────────────────────────────────
\section{Criação das Páginas HTML}

\subsection{Criando a estrutura de diretórios}

\begin{lstlisting}[style=bash]
# Exemplo do diretorio de um site.
mkdir -p /var/www/marataizes.joao.info
\end{lstlisting}

\subsection{Criando as páginas HTML (conteúdo distinto para cada site)}

Cada página deve identificar claramente o host e o domínio ao qual pertence.
Exemplo para \texttt{marataizes.joao.info}:

\begin{lstlisting}[style=html]
<!DOCTYPE html>
<html lang="pt-BR">
<head>
  <meta charset="UTF-8">
  <title>Marataizes - joao.info</title>
  <style>body { font-family: sans-serif; text-align: center; margin-top: 60px; }</style>
</head>
<body>
  <h1>Marataizes</h1>
  <p>Dominio: <strong>joao.info</strong></p>
  <p>Cidade do sul do Espirito Santo, conhecida pela producao de abacaxi.</p>
</body>
</html>
\end{lstlisting}

Crie os seis arquivos com conteúdo diferente:

\begin{lstlisting}[style=bash]
vim /var/www/marataizes.joao.info/index.html
vim /var/www/castelo.joao.info/index.html
vim /var/www/cachoeiro.joao.info/index.html
vim /var/www/marataizes.maria.info/index.html
vim /var/www/castelo.maria.info/index.html
vim /var/www/cachoeiro.maria.info/index.html
\end{lstlisting}

\begin{alertbox}
  \textbf{Atenção:} Cada página deve ter um texto diferente das demais. Usar o mesmo conteúdo em páginas distintas implica perda de pontuação no critério de páginas HTML.
\end{alertbox}

% ── Seção 4 ───────────────────────────────────────────────────────────────────
\section{Configuração dos Virtual Hosts no NGINX}

\subsection{Criando os arquivos de configuração}

Crie um arquivo em \texttt{/etc/nginx/sites-available/} para cada site.
Modelo para \texttt{marataizes.joao.info}:

\begin{lstlisting}[style=conf]
# /etc/nginx/sites-available/marataizes.joao.info
server {
    listen 80;
    server_name marataizes.joao.info;

    root /var/www/marataizes.joao.info;
    index index.html;

    location / {
        try_files $uri $uri/ =404;
    }
}
\end{lstlisting}

Repita para os outros cinco sites, alterando \texttt{server\_name} e \texttt{root}:

\begin{lstlisting}[style=bash]
vim /etc/nginx/sites-available/castelo.joao.info
vim /etc/nginx/sites-available/cachoeiro.joao.info
vim /etc/nginx/sites-available/marataizes.maria.info
vim /etc/nginx/sites-available/castelo.maria.info
vim /etc/nginx/sites-available/cachoeiro.maria.info
\end{lstlisting}

\subsection{Ativando os sites}

Crie os links simbólicos para habilitar cada site:

\begin{lstlisting}[style=bash]
    ln -s /etc/nginx/sites-available/marataizes.maria.info \
          /etc/nginx/sites-enabled/marataizes.maria.info
\end{lstlisting}

\subsection{Testando a configuração e recarregando o NGINX}

\begin{lstlisting}[style=bash]
nginx -t
systemctl reload nginx
systemctl status nginx
\end{lstlisting}

\begin{successbox}
  A saída de \texttt{nginx -t} deve informar \texttt{syntax is ok} e \texttt{test is successful}. Qualquer erro indica problema em um dos arquivos de configuração.
\end{successbox}

% ── Seção 5 ───────────────────────────────────────────────────────────────────
\section{Testes e Verificação}

\subsection{Configurando o DNS no cliente}

Na máquina cliente, configure o servidor DNS para apontar para \texttt{<IP\_VM>}.
No Debian/Ubuntu, edite o arquivo de resolução:

\begin{lstlisting}[style=bash]
# Edite /etc/resolv.conf
vim /etc/resolv.conf

# Adicione (ou substitua a linha nameserver existente):
nameserver <IP\_DNS>
\end{lstlisting}

\begin{infobox}
  Em ambientes com NetworkManager, a configuração do DNS deve ser feita pela interface gráfica de rede ou pelo comando \texttt{nmcli con mod <conexão> ipv4.dns <IP\_VM>}.
\end{infobox}

\subsection{Testes de resolução DNS}

Execute no cliente os comandos abaixo e registre as saídas no tutorial:

\begin{lstlisting}[style=bash]
# Verificar resolucao dos seis nomes
dig marataizes.joao.info
dig castelo.joao.info
dig cachoeiro.joao.info
dig marataizes.maria.info
dig castelo.maria.info
dig cachoeiro.maria.info
\end{lstlisting}

A seção \texttt{ANSWER SECTION} de cada resultado deve mostrar o IP \texttt{<IP\_VM>}.
Exemplo de saída esperada:

\begin{lstlisting}[style=conf]
;; ANSWER SECTION:
marataizes.joao.info.  38400  IN  A  <IP\_VM>
\end{lstlisting}

\subsection{Testes de acesso HTTP}

\begin{lstlisting}[style=bash]
Utilize no navegador: http://marataizes.joao.info
\end{lstlisting}

O conteúdo HTML de cada site deve ser exibido no terminal. Acesse também pelo \textbf{navegador} do cliente para confirmar visualmente que cada URL exibe uma página diferente.

% ── Seção 6 ───────────────────────────────────────────────────────────────────
\section{O que o Tutorial (PDF) deve conter}

O documento a ser entregue deve apresentar as seguintes evidências:

\begin{enumerate}
  \item \textbf{Identificação:} nomes dos integrantes da dupla e os domínios utilizados.
  \item \textbf{Configuração do DNS via Webmin:}
    \begin{itemize}
      \item Capturas de tela das zonas criadas no BIND DNS Server.
      \item Capturas de tela dos registros A de cada zona.
      \item Conteúdo dos arquivos \texttt{/etc/bind/db.nome1.info} e \texttt{/etc/bind/db.nome2.info}.
      \item Saída de \texttt{systemctl status bind9} mostrando o serviço ativo.
    \end{itemize}
  \item \textbf{Configuração do NGINX:}
    \begin{itemize}
      \item Conteúdo dos seis arquivos de virtual host em \texttt{sites-available/}.
      \item Saída do comando \texttt{nginx -t} e \texttt{systemctl status nginx}.
    \end{itemize}
  \item \textbf{Páginas HTML:} conteúdo de cada um dos seis arquivos \texttt{index.html}.
  \item \textbf{Testes:}
    \begin{itemize}
      \item Saída do \texttt{dig} para os seis nomes de domínio.
      \item Saída do \texttt{curl} \textbf{ou} capturas de tela do navegador exibindo os seis sites com conteúdos distintos.
    \end{itemize}
\end{enumerate}

% ── Critérios de avaliação ─────────────────────────────────────────────────────
\vspace{0.5em}
\begin{center}
  \Large\textbf{Critérios de Avaliação}
\end{center}

\begin{center}
  \begin{tabular}{|p{10.5cm}|c|}
    \hline
    \textbf{Critério} & \textbf{Pontuação} \\
    \hline
    Zonas DNS criadas corretamente no BIND via Webmin (2 zonas) & 2,0 \\
    \hline
    Registros A configurados para os 6 hosts (3 por domínio) & 2,0 \\
    \hline
    Virtual hosts do NGINX configurados corretamente para os 6 sites & 2,0 \\
    \hline
    Páginas HTML com conteúdo distinto e identificável para cada site & 2,0 \\
    \hline
    Testes de resolução DNS documentados (\texttt{dig}/\texttt{nslookup}) & 1,0 \\
    \hline
    Testes de acesso HTTP documentados (\texttt{curl}/navegador) & 1,0 \\
    \hline
    \textbf{Total} & \textbf{10,0} \\
    \hline
  \end{tabular}
\end{center}

\vspace{0.5em}

\begin{alertbox}
  \textbf{Observação:} A dupla deve usar obrigatoriamente os primeiros nomes dos dois integrantes como nomes de domínio. O tutorial entregue deve conter evidências claras (capturas de tela e/ou saídas de terminal) de \textbf{todas} as etapas realizadas. Trabalhos sem evidências de testes não serão pontuados nos critérios 5 e 6.
\end{alertbox}

\end{document}
